\documentclass[10pt,a4paper]{article}
\usepackage[margin=1.8cm]{geometry}
\usepackage[T1]{fontenc}
\usepackage{amsmath}
\usepackage[version=4]{mhchem}
\usepackage{enumitem}
\usepackage{titlesec}
\usepackage{xcolor}
\usepackage{tcolorbox}

\setlist{nosep,leftmargin=1.4em}
\titleformat{\section}{\large\bfseries}{}{0em}{}
\titlespacing*{\section}{0pt}{1.1ex}{0.5ex}
\setlength{\parindent}{0pt}
\setlength{\parskip}{0.4ex}
\pagestyle{empty}

\begin{document}

{\Large\bfseries LRP 241: Meeting 2 Plan}\\[0.3ex]
{\small GWP of Hydrogen \quad|\quad Supervisor: Yannis Hardalupas \quad|\quad Papers: Sand et al.\ 2023; Hauglustaine et al.\ 2022}

\section{Purpose}
\begin{enumerate}
  \item Show you have read and understood the two assigned papers.
  \item Get scope decisions only he can make.
  \item Leave with agreed next steps recorded in the meeting log.
\end{enumerate}
He is testing whether you critique or merely summarise.

\section{Running order (25--30 min)}
\begin{enumerate}[label=\textbf{\arabic*.}]
  \item \textbf{Since meeting 1 (1 min).} 26 sources tiered; tracker built; Krishnan et al.\ 2026 located. Facts only.

  \item \textbf{Headline (30 s). Lead with this.}
  \begin{tcolorbox}[colback=gray!8,colframe=gray!50,boxsep=2pt,left=4pt,right=4pt,top=2pt,bottom=2pt]
  \emph{``The apparent convergence on a GWP of $\approx 12$ is models sharing assumptions. The real uncertainty reduces to two physical quantities: how much leaked \ce{H2} the soil absorbs, and where it leaks.''}
  \end{tcolorbox}

  \item \textbf{Sand 2023 (3 min).}
  \begin{itemize}
    \item 5-model ensemble; $+10\%$ surface \ce{H2}; steady state $\Rightarrow \mathrm{GWP}_{100} = 11.6 \pm 2.8$.
    \item Settles Derwent's pulse vs steady-state objection.
    \item Does not settle: soil-sink uncertainty is expert judgement and the largest term; UKCA omits stratospheric \ce{H2O}; emission location untested (authors admit, p7).
  \end{itemize}

  \item \textbf{Hauglustaine 2022 (3 min).}
  \begin{itemize}
    \item One model plus analytical impulse-response functions $\Rightarrow \mathrm{GWP}_{100} = 12.8 \pm 5.2$.
    \item Key point: their own NH-only value is $10.7$, which supports Derwent's latitude argument.
    \item Break-even thresholds (feed objective 4): green \ce{H2} beneficial at any realistic leak; $>60\%$ blue and $>7\%$ leak $\Rightarrow$ no benefit at 100 yr.
    \item Derwent's borrowed stratospheric term: $12.8 \times 0.47 \times 0.61 \approx 3.7$, consistent.
  \end{itemize}

  \item \textbf{Tie the four papers together (2 min).} Read the ``Comparison'' section of the notes. This is where you show critique.

  \item \textbf{Objective 4 started (2 min).} Hauglustaine (thresholds); Frazer-Nash 2022 (stage-by-stage leak rates); EA/NPL 2025 (detection is built for safety at $\sim 4\%$ LFL, not for climate-relevant chronic leaks $<1\%$).

  \item \textbf{Report plan (1 min).} Page split and figures.

  \item \textbf{Questions. Most important; cut items 3--7 before cutting this.}
  \begin{enumerate}[label=Q\arabic*.]
    \item Page split OK? Objective 4 capped at $\sim 3.5$ pages?
    \item \textbf{Highest priority.} Is a simple break-even calculation acceptable as own analysis:
    \[
      \mathrm{GWP}_{100}\, f_{\mathrm{leak}}\, m_{\ce{H2}} \quad \text{vs} \quad m_{\ce{CO2},\,\mathrm{avoided}} \;?
    \]
    \item Include Krishnan 2026 as the current state of the soil-sink question?
    \item Cite BEIS ($11 \pm 5$), ACP ($12 \pm 6$), or both?
    \item Redraw published figures or reproduce with attribution?
    \item Meeting 3 date.
  \end{enumerate}
  \textbf{Extra:} Hauglustaine p10 states an OH change of ``$+8\%$'' per 200\,Tg/yr; added \ce{H2} should deplete OH. Ask whether you are misreading it. Frame as a question.

  \item \textbf{AI reflection (1 min).} Used for screening, extraction and challenging analysis; every number checked against the PDF; three errors caught; no report text.

  \item \textbf{Close.} Agree: Tier B extraction; section 1 draft and section 2 skeleton before meeting 3; next papers Krishnan 2026 and Ocko \& Hamburg 2022. Repeat them back to him.
\end{enumerate}

\section{Know these cold (he will probe)}
\begin{itemize}
  \item \textbf{Soil sink:} $\sim 75\%$ of \ce{H2} removal. Larger sink $\Rightarrow$ shorter lifetime $\Rightarrow$ lower GWP.
  \item \textbf{\ce{CH4} feedback factor:} less OH $\Rightarrow$ longer \ce{CH4} lifetime $\Rightarrow$ even less OH; amplifies the \ce{CH4} term by $\approx 1.3$--$1.5$.
  \item \textbf{Steady state vs pulse:} the time-integrated response to a pulse equals the steady-state perturbation scaled by lifetime, so the two are equivalent.
  \item \textbf{Why NH is lower:} more land $\Rightarrow$ more soil uptake $\Rightarrow$ shorter lifetime ($2.0$ vs $2.5$\,yr) $\Rightarrow$ less warming.
\end{itemize}

\section{After the meeting}
\begin{enumerate}
  \item Read the meeting log before clicking Agree; it records what he approved, especially Q2.
  \item Same-day follow-up email: agreed actions and meeting 3 date.
  \item Log his answers in the tracker.
\end{enumerate}

\end{document}
